\scanpage{77}
\mspage{34}
\uline{15}. Introduction des proespaces :
${\unsure{V|\mathcal V}}_{i_0.\!-.i_n},\ {\unsure{V|\mathcal V}}_J^{\deleted{\unsure{*}}L}$,
% The title has a dot, a short horizontal dash, and a dot between the indices.
% Both V forms have a curved bowl; plain/calligraphic alphabet assignment remains open.

On désigne par ${\unsure{V|\mathcal V}}_i^{\unsure{*|<|\deleted{*}}}$
\unsure{\deleted{le}|le} \inserted{pro}voisinage tubulaire
(au sens des proespaces) de $X_i^*=X_i\setminus\bigcup_{j<i}X_j$
dans $X$. Ainsi, \deleted{\unsure{on}} \unsure{on a}
% The raised pro has a curved leader ending immediately before voisinage.
\[
(51)\qquad X_i^*\longrightarrow {\unsure{V|\mathcal V}}_i^{\unsure{*|\deleted{*}}}
\qquad\begin{gathered}
\text{équivalence d'homo-}\\
\text{topie \unsure{des|de} proespaces.}
\end{gathered}
\]
% The hyphen after homo is written at the right paper edge.

\deleted{\unsure{Pour un drapeau}}
Si\unsure{)} $i<j$, on pose
% The separate curved stroke after Si is retained; no opening parenthesis is supplied.
\[
(52)\qquad
\unsure{\rlap{$X$}V|\rlap{$X$}\mathcal V}_{ij}={\unsure{V|\mathcal V}}_i\cap {\unsure{V|\mathcal V}}_j
\]
% The left base in (52) is one superposed X/V cluster with a single shared ij.
% The overprinting order and deletion state cannot be separated; no second index is supplied.

\margin{Pour un
\diagram{\begin{tikzcd}[column sep=small]
{\unsure{V|\mathcal V}}_{ij}^{\unsure{b}}={\unsure{V|\mathcal V}}_i\cap X_j^*\arrow[r,hook]&{}
\end{tikzcd}}
est-ce une équivalence d'homotopie ?}
% The oblique left note ends its formula at the curved arrowhead; no target letters are visible.

\margin{NB ${\unsure{V|\mathcal V}}_{ij}$ a un sens \deleted{\unsure{si}} si $i\neq j$
mais si $i$ et $j$ ne sont pas comparables, on a
${\unsure{V|\mathcal V}}_i\cap {\unsure{V|\mathcal V}}_j=\varnothing$,
\inserted{(du moins si les ouverts de $X$ sont normaux, p.ex. $X$ métrisable)}
si $i=j$, on a ${\unsure{V|\mathcal V}}_{ii}={\unsure{V|\mathcal V}}_i$}
% The long curved leader joins the parenthetical insertion to the intersection-empty clause.
% The parenthetical words are physically below V_ii=V_i; their attachment precedes si i=j.

de sorte qu'on a
\diagram{\begin{tikzcd}[column sep=large,row sep=huge]
&{\unsure{V|\mathcal V}}_{ij}\arrow[r,hook,"b_{ij}"]\arrow[d,hook,"s_{ij}"]&{\unsure{V|\mathcal V}}_j\\
\Bigl(X_i
\arrow[r,shift right=1.2ex]
\arrow[r,phantom,"{\substack{\unsure{(w)|(h)}\\\sim\\
\unsure{\longleftarrow|\deleted{\longleftarrow}}}}"]
&\Bigr){\unsure{V|\mathcal V}}_i&
\end{tikzcd}}
% Four written connectors: two hooked main arrows, the dense left-pointing lower mark,
% and a separate lower right-pointing arrow. The left-pointing mark has unresolved strike state.
% The closing parenthesis lies before V_i, with the one tilde and (w)/(h) above the lower marks.
On a :
\[
(53)\qquad X=\varinjlim\bigl({\unsure{V|\mathcal V}}_i,{\unsure{V|\mathcal V}}_{ij},
\unsure{\deleted{\}}|\}}\ b_{ij},s_{ij}\bigr)
\]
% The brace-like glyph between the comma and b is crossed/reworked; its strike state remains open.
$X$ est obtenu par \uline{recollement} des ${\unsure{V|\mathcal V}}_i$
suivant les ${\unsure{V|\mathcal V}}_{ij}$.

On définit aussi, plus généralement, les
\[
(54)\qquad {\unsure{V|\mathcal V}}_{i_0.\!-i_n}={\unsure{V|\mathcal V}}_{i_0}\cap
\mathord{-}\,\mathord{-}\cap {\unsure{V|\mathcal V}}_{i_n}
\]
% The left index has dot-dash; the right continuation has two separate dashes.
% The final V is retraced, without a recoverable independently cancelled underlying letter.
pour un drapeau $i_0<i_1<.\!-<i_n$
\margin{\uline{NB} L'expression a un sens $=$ si
l'ens. $\{i_0,\mathord{-}i_n\}$ ne forme pas un drapeau,
mais il est vide si ce n'est pas un drapeau}
% The bottom marginal note is oblique. Its = has two horizontal strokes, the lower slightly wavy.
% The continuation within the written braces is one dash, with no second comma before i_n.
% Pale bottom-left foliation in another, unnamed hand:
\margin{76}
